\section{Pioneering, conformity, and the trademark record}\label{sec:related}

\subsection{Pioneering and optimal distinctiveness}

Whether market pioneers are advantaged or burdened is an old, unresolved
question measured mostly on named markets and survivor-prone samples
\citep{GolderTellis1993, SemadeniAnderson2010}. A parallel literature
finds evaluators reward conformity with a penalty for extremity in both
directions --- the optimal distinctiveness of \citet{Zuckerman1999},
\citet{Hsu2006}, \citet{Zhao2017}, and \citet{TaeuscherRothe2020}. This
paper meets both at census scale: every entrant's language is observed,
whether or not the product survived to be remembered, and the record
supplies two different evaluators --- an examiner at registration and the
market at the five-year proof of continued use --- who turn out to price
the two components differently. The examiner rewards middling conformity
and cannot observe lead; the market is indifferent to conformity once
theme importation is held, and penalizes having been early.

\subsection{Trademarks as innovation data}

A trademark filing is a dated, costly, public disclosure made by the
millions of firms that never patent, and a literature has grown around
reading it as an innovation indicator \citep{Castaldi2020, Willeke2026}.
Trademark counts and breadth predict venture valuations
\citep{Block2014}; filing timing tracks innovation processes
\citep{Seip2018}; and trademarks mark the product-market side of
innovation that patents miss \citep{Flikkema2019, Flikkema2025}.
Census-linked work reads registration as a correlate of firm growth and
survival \citep{Dinlersoz2018, Dinlersoz2023}. Nearly all of it counts
filings, classes, or renewals; the descriptions' text enters, if at all,
as breadth or similarity. This paper reads the text itself, and reads it
against time.

\subsection{Text measures of novelty}

Scoring a document by its divergence from a reference corpus before and
after it comes from cultural analytics: \citet{Murdock2017} on Darwin's
reading notebooks, \citet{Barron2018} on the French Revolution's
parliamentary debates. In economics and finance, text-based measures
locate firms in product space \citep{HobergPhillips2016}, date
technological innovation in patents over the long run \citep{Kelly2021},
and track the vintage of patent language \citep{Kalyani2024}. The
construction here differs in its object --- the composition of commercial
offerings, not ideas or claims --- and in its two-sided reference: every
filing is scored against its own industry's language both before and
after its date, which is what lets \emph{unusual} be separated from
\emph{early}. The separation matters because the two components predict
opposite things.

\subsection{Nomenclature}\label{sec:nomenclature}

The innovation literature is not consistent about what quantities like
these are called, and several of its established terms assert things this
design cannot show. \emph{Novelty} and \emph{transience}
\citep{Murdock2017} name the two KL levels themselves, and novelty is
declined because it names only the backward-looking half of a quantity
that faces both ways; \emph{resonance} \citep{Barron2018} is
arithmetically what this paper calls lead and is avoided because it names
a mechanism the design cannot observe; \emph{radicalness},
\emph{originality} and \emph{generality} \citep{HallJaffe2001,
Verhoeven2025} each claim something about the product or its realized
influence that a description's position cannot carry. The companion paper
gives the full argument for each term; Table~\ref{tab:nomenclature} maps
the five quantities to their nearest occupants and to the terms used
here.


\begin{table}[h]\centering\small
\caption{The five quantities, their nearest occupants in the literature,
and the terms used here. $P$ is a filing's theme distribution;
the references $Q^{-}$ and $Q^{+}$ aggregate the five years of
same-class filings before and after it (Section~\ref{sec:measure}).
\emph{Lagging} and \emph{leading} describe the two signs of $L$. The
level names appear only where the two-sided construction itself is at
issue; the analysis runs on $A$ and $L$.}
\label{tab:nomenclature}
\begin{tabular}{p{3.4cm}p{5.6cm}p{4.4cm}}
\toprule
Quantity & Nearest in the literature & Term used here \\
\midrule
$K^{-} = \mathrm{KL}(P \| Q^{-})$ & novelty \citep{Murdock2017, Barron2018}; originality \citep{HallJaffe2001} & past-facing surprise \\
$K^{+} = \mathrm{KL}(P \| Q^{+})$ & transience \citep{Murdock2017}; generality \citep{HallJaffe2001}, which measures realized influence, not position & future-facing surprise \\
$A = \tfrac{1}{2}(K^{-} + K^{+})$ & atypicality \citep{Uzzi2013}, which scores internal pairings, not distance from the class & \textbf{atypicality} \\
$L = K^{-} - K^{+}$ & resonance \citep{Barron2018}, which asserts a mechanism & \textbf{lead} (leading / lagging) \\
$|L|$ & no established occupant & lead magnitude \\
\bottomrule
\end{tabular}
\end{table}
